\SourceLeaf{220}
\Needspace{12\baselineskip}
\begin{center}
\textit{Allgemeine Lösung aller unbestimmten Gleichungen zweiten Grades mit zwei Unbekannten in ganzen Zahlen.}
\end{center}

Zum Abschluss dieser Untersuchungen fügen wir hinzu:

\Needspace{8\baselineskip}
\Article{216}
\textsc{Aufgabe.} \textit{Alle ganzzahligen Lösungen der allgemeinen}\HistoricalNote{*)}{Ist eine Gleichung gegeben, deren zweiter, vierter oder fünfter Koeffizient nicht gerade ist, so erhält sie durch Multiplikation mit 2 die hier vorausgesetzte Form.} \textit{unbestimmten Gleichung zweiten Grades mit zwei Unbekannten}
\[
axx+2bxy+cyy+2dx+2ey+f=0
\]
\textit{finden (wobei $a,b,c$ usw. beliebige gegebene ganze Zahlen sind).}

\textit{Lösung.} Anstelle der Unbekannten $x,y$ führen wir
\[
p=(bb-ac)x+be-cd\quad\text{und}\quad q=(bb-ac)y+bd-ae
\]
ein. Diese sind offenbar immer ganze Zahlen, wenn $x,y$ ganze Zahlen sind. Damit erhalten wir die Gleichung
\[
app+2bpq+cqq+f(bb-ac)^2+(bb-ac)(aee-2bde+cdd)=0
\]
oder, wenn wir der Kürze wegen
\[
f(bb-ac)^2+(bb-ac)(aee-2bde+cdd)=-M
\]
setzen, die Gleichung
\[
app+2bpq+cqq=M
\]
Wir haben im Vorhergehenden gezeigt, wie wir alle Lösungen dieser Gleichung, \textit{d.h.} alle Darstellungen von $M$ durch $(a,b,c)$, finden können. Bestimmen wir nun aus den einzelnen Werten von $p,q$ die entsprechenden Werte von $x,y$ mithilfe der Gleichungen
\[
x=\frac{p+cd-be}{bb-ac},\quad y=\frac{q+ae-bd}{bb-ac}
\]
so sehen wir leicht, dass alle diese Werte die gegebene Gleichung erfüllen und dass es keine ganzzahligen Werte von $x,y$ gibt, die wir auf diese Weise nicht erhalten. Verwerfen wir also unter allen so entstehenden Werten von $x,y$ die gebrochenen, so bleiben alle gesuchten Lösungen übrig.

Zu dieser Lösung ist Folgendes zu bemerken.

\SourceLeaf{221}
$1^{\mathrm{o}}$ Lässt sich entweder $M$ nicht durch $(a,b,c)$ darstellen, oder ergeben sich aus keiner Darstellung \textit{ganzzahlige} Werte von $x,y$, so hat die Gleichung keine ganzzahlige Lösung.

$2^{\mathrm{o}}$ Ist die Determinante von $(a,b,c)$, \textit{d.h.} $bb-ac$, negativ, oder ist sie positiv quadratisch und ist zugleich $M$ nicht $=0$, so ist die Anzahl der Darstellungen von $M$ durch $(a,b,c)$ endlich. Daher ist auch die Anzahl aller Lösungen der gegebenen Gleichung, wenn es überhaupt welche gibt, endlich.

$3^{\mathrm{o}}$ Ist $bb-ac$ positiv und nichtquadratisch oder quadratisch und zugleich $M=0$, so lässt sich $M$, wenn überhaupt, auf \textit{unendlich viele verschiedene Arten} durch $(a,b,c)$ darstellen. Da es aber unmöglich ist, alle diese Darstellungen \textit{selbst} zu finden und zu prüfen, ob sie ganzzahlige oder gebrochene Werte von $x,y$ liefern, brauchen wir eine Regel, mit der wir \textit{Gewissheit} erlangen können, falls überhaupt keine Darstellung ganzzahlige Werte von $x,y$ liefert (denn wie viele Darstellungen wir in diesem Fall auch \textit{prüfen}, ohne eine solche Regel würden wir nie Gewissheit erreichen). Liefern dagegen einige Darstellungen ganzzahlige Werte von $x,y$, andere gebrochene, so müssen wir zeigen, wie wir die letzteren von den ersteren im Voraus allgemein unterscheiden können.

$4^{\mathrm{o}}$ Ist $bb-ac=0$, so lassen sich die Werte von $x,y$ mit den vorhergehenden Formeln überhaupt nicht bestimmen. Daher müssen wir für diesen Fall eine \textit{besondere Methode} finden.

